\section{The exact hidden-form experiment in its critical window}

Fix a prime power $q$, let $n\ge4$ tend to infinity through even integers, and
put $D=\binom n2$ and $d=D-1$. Let $\mathcal H_n$ be the projective
nondegenerate alternating forms on $\mathbb F_q^n$, and let
$\mathcal G_n=\operatorname{Gr}(2,\mathbb F_q^n)$. Under the null law
$Q_{n,t}$ the observed planes are independent and uniform on $\mathcal G_n$.
Under the planted law $P_{n,t}$ a uniform $[H]\in\mathcal H_n$ is chosen
once, and the planes are conditionally independent uniform $H$-isotropic
planes. Define the compatible linear space
\[
 \mathcal N=\{G\in\operatorname{Alt}_n(q):G|_{K_i\times K_i}=0
                 \text{ for all observed }K_i\},
\]
and let $Z_n$ be the number of projective nondegenerate forms in
$\mathcal N$. We use
\[
 c_q=\prod_{j\ge2}(1-q^{-(2j-1)})=\frac{\beta_q}{1-q^{-1}},
 \qquad g_j=\prod_{i=1}^j(1-q^{-i}),\qquad g_0=1.
\]

For each integer $b$, define a probability distribution on
$k\ge\max(0,-b)$ by
\[
 \pi_{q,b}(k)=\frac{g_\infty q^{-k(k+b)}}{g_kg_{k+b}}.
\]
This is the limiting nullity law of \cref{cor:nullity}, with $b=u$. Write
$m_k=(q^k-1)/(q-1)$.

\begin{theorem}[Complete critical law and optimal inference]
\label{thm:complete-critical-law}
Fix an integer $a$, put $t=d+a$ and $\lambda=c_qq^{-a}$, and define a
limiting pair $(K,Z)$ by
\[
 \Pr(K=k)=\pi_{q,a-1}(k),\qquad
 (Z\mid K=k)\sim\operatorname{Bin}(m_k,\beta_q).
\]
Then under $Q_{n,d+a}$,
\[
 (\dim\mathcal N,Z_n)\longrightarrow(K,Z)
\]
in total variation. Under the planted law, the limiting pair instead has
\[
 \Pr(K=k)=\pi_{q,a}(k-1),\qquad
 (Z\mid K=k)\sim1+\operatorname{Bin}(m_k-1,\beta_q),
 \qquad k\ge1.
\]
The exact limiting full-observation total variation is
\begin{align*}
 T_q(a)
 &=\lim_n\operatorname{TV}(P_{n,d+a},Q_{n,d+a})\\
 &=\mathbb E(1-Z/\lambda)_+
  =\frac12\mathbb E|Z/\lambda-1|\\
 &=\sum_{k\ge\max(0,1-a)}\pi_{q,a-1}(k)
   \sum_{z=0}^{m_k}\binom{m_k}{z}\beta_q^z(1-\beta_q)^{m_k-z}
                         (1-z/\lambda)_+.
\end{align*}
For $a\ge0$ this reduces to the positive series
\[
 T_q(a)=\sum_k\pi_{q,a-1}(k)(1-\beta_q)^{m_k}.
\]
The limiting optimal probability of recovering the hidden projective form is
\[
 S_q(a)=\frac{1-\sum_k\pi_{q,a-1}(k)(1-\beta_q)^{m_k}}{\lambda}.
\]
The limiting probability that there is a unique compatible nondegenerate
projective form is
\[
 U_q(a)=\frac{\beta_q}{\lambda}
     \sum_{k\ge1}\pi_{q,a-1}(k)m_k(1-\beta_q)^{m_k-1}.
\]
At each finite size the optimal equal-prior test compares $Z_n$ to
$\lambda_n=|\mathcal H_n|p_n^t$, where
\[
 p_n=\frac{q^{n-2}-1}{q^{n-1}-1}.
\]
Thus counting the compatible nondegenerate forms retains all information
for testing. For $a\ge0$, accepting the planted model exactly when
$Z_n>0$ is already optimal at every $n\ge4$.
\end{theorem}

\begin{proof}
The exact group-orbit and plane counts give
\[
 |\mathcal H_n|=q^{D-1}c_n,\qquad
 c_n=\prod_{j=2}^{n/2}(1-q^{-(2j-1)}),\qquad
 \frac{dP_{n,t}}{dQ_{n,t}}=\frac{Z_n}{\lambda_n}.
\]
The first identity follows by dividing the number of nondegenerate
alternating forms by $q-1$. For a fixed $H$, count ordered independent
vectors $(u,v)$ with $H(u,v)=0$ to obtain $p_n$; this probability does not
depend on $H$. Each compatible projective form contributes the same
conditional density $p_n^{-t}$, proving the likelihood formula.
Since $p_n=q^{-1}(1+O_q(q^{2-n}))$ and $t=O(n^2)$,
$\lambda_n\to c_qq^{-a}=\lambda$.

Theorem~\ref{thm:planes} makes the law of $\mathcal N$ under $Q$
asymptotically the law of the kernel of a uniform $t\times D$ matrix.
Conditional on its dimension, that latter kernel is uniform among
subspaces of that dimension. Its dimension tends to $\pi_{q,a-1}$.
Corollary~\ref{cor:binomial-pencil} gives the conditional binomial limit
for each fixed dimension. Truncating the tight dimension distribution
proves total-variation convergence of the pair.

For this null limit,
\[
 \mathbb E q^K=1+q^{1-a},\qquad
 \mathbb E m_K=\frac{q^{1-a}}{q-1},\qquad
 \mathbb EZ=\beta_q\mathbb E m_K=c_qq^{-a}=\lambda.
\]
The first identity follows either by summing $\pi_{q,a-1}$ or by counting
nonzero kernel vectors in the uniform rectangular matrix model.
The planted limit has point masses $z/\lambda$ times the null point
masses. They sum to one by the last identity. Pointwise convergence of
the exact likelihood-weighted masses, followed by preservation of total
mass, gives total-variation convergence under the planted law too.
Size-biasing a binomial variable gives $1+\operatorname{Bin}(m_k-1,\beta_q)$.
The size-biased dimension weight is
\[
 \pi_{q,a-1}(k)\frac{\beta_qm_k}{\lambda}
 =\pi_{q,a-1}(k)q^{a-1}(q^k-1)=\pi_{q,a}(k-1).
\]

For any likelihood $L$, total variation equals
$\mathbb E_Q(1-L)_+$. This is a bounded continuous function of $L$, so
the convergence just proved yields $T_q(a)$ directly. Since
$\mathbb EZ=\lambda$, the absolute-value form follows as well. If $a\ge0$
then $0<\lambda<1$, so only $Z=0$ contributes. At finite size one has
$p_n<q^{-1}$ and $c_n<1$ for $n\ge4$, whence
$\lambda_n<c_nq^{-a}<1$ and the same testing rule is exactly optimal.

The posterior distribution of $[H]$ is uniform over the $Z_n$ compatible
nondegenerate forms. Hence the optimal recovery success is
\[
 \mathbb E_P(1/Z_n)
 =Q_{n,t}(Z_n>0)/\lambda_n.
\]
This proves the recovery formula. Similarly,
$P_{n,t}(Z_n=1)=Q_{n,t}(Z_n=1)/\lambda_n$, giving the uniqueness formula.
\end{proof}

\paragraph{Arbitrary prior probabilities.}
If the prior probability of the planted law is $\theta\in(0,1)$, its
optimal limiting Bayes error is
\[
 \mathbb E\min\{1-\theta,\theta Z/\lambda\}.
\]
At finite size the optimal test accepts the planted model when
$Z_n/\lambda_n>(1-\theta)/\theta$, with either decision allowed at a tie.
This follows from the same exact likelihood and requires no further
limit theorem.

\paragraph{A rigorous truncation bound.}
For $r\ge\max(1,1-a)$,
\[
 \Pr(K\ge r)\le\mathbb E\qbinom{K}{r}
 =\frac{q^{r(1-a)-r^2}}{g_r}.
\]
The identity follows from counting the $r$-spaces in a uniform matrix
kernel before taking its dimension to infinity. Thus truncating the
positive series for $T_q(a)$ before $k=r$ incurs error no larger than the
displayed bound; the same statement holds for the bounded-summand formula
when $a<0$. Product tails may be enclosed by
$1-\sum_{j>J}q^{-j}\le\prod_{j>J}(1-q^{-j})\le1$, with the analogous
odd-index bound for $\beta_q$.

\begin{corollary}[Exact distinction from dimension-only inference]
\label{cor:dimension-gap}
The limiting total variation when only $\dim\mathcal N$ is retained is
\[
 T_q^{\rm dim}(a)=
 \begin{cases}
 q^{a-1}g_\infty/g_{1-a},&a\le0,\\
 g_\infty/g_{a-1},&a\ge1.
 \end{cases}
\]
For every fixed integer $a$, the full-observation distance is strictly
larger than the dimension-only distance. As $a\to-\infty$ after the
critical-window limit,
\[
 T_q(a)\sim\sqrt{\frac{1-\beta_q}{2\pi c_q}}\,q^{a/2},
 \qquad T_q^{\rm dim}(a)\sim q^{a-1}.
\]
\end{corollary}

\begin{proof}
The dimension likelihood is
$\mathbb E(Z/\lambda\mid K=k)=q^{a-1}(q^k-1)$.
For $a\le0$ put $h=1-a\ge1$. The smallest possible null dimension is $h$;
its likelihood is $1-q^{-h}<1$, whereas all larger dimensions have
likelihood greater than one. Thus the distance is
$q^{-h}\pi_{q,-h}(h)=q^{a-1}g_\infty/g_{1-a}$.
For $a\ge1$, only $k=0$ has likelihood below one, except the harmless tie
at $q=2,a=1,k=1$. This gives $\pi_{q,a-1}(0)=g_\infty/g_{a-1}$.

For strictness, choose any sufficiently large $k$ with positive
$\pi_{q,a-1}(k)$ and $m_k>\lambda$. Conditional on that dimension, the
binomial variable has positive probability both at $0$ and above
$\lambda$. The function $(1-x/\lambda)_+$ is therefore strictly above
its Jensen value after averaging over that conditional distribution.
Summing gives a strict gap.

For the asymptotic formula put $h=1-a\to\infty$. Then
\[
 \Pr(K=h)=g_\infty/g_h=1-O_q(q^{-h}),\qquad
 \lambda=\frac{\beta_qq^h}{q-1}.
\]
On $K=h$, $Z$ is binomial with $m_h=(q^h-1)/(q-1)$ trials, mean
$\lambda-\beta_q/(q-1)$ and variance asymptotic to
$(1-\beta_q)\lambda$. The usual binomial central limit theorem, together
with bounded second moments of its normalized version, gives
\[
 \mathbb E|Z-\lambda|\sim
 \sqrt{2/\pi}\sqrt{(1-\beta_q)\lambda}
\]
for this conditional binomial law. The omitted dimensions contribute
only $O_q(\lambda^{-1})$ to the normalized absolute-value expectation:
indeed
\[
 \mathbb E[m_K;K>h]
 =\frac{q^h}{q-1}-\frac{q^h-1}{q-1}\frac{g_\infty}{g_h}=O_q(1),
\]
and $\Pr(K>h)=O_q(q^{-h})=O_q(\lambda^{-1})$.
Dividing by $2\lambda$ proves the first asymptotic; the exact dimension
formula gives the second.
\end{proof}

\paragraph{Optimal inference with bounded expected enumeration.}
Compute a basis of $\mathcal N$ by Gaussian elimination, enumerate its
projective points, and test the corresponding alternating matrices for
nonsingularity. This computes $Z_n$ and hence the optimal test and the
uniform posterior exactly. At every fixed critical offset, the expected
number of enumerated points stays bounded under both hypotheses.
Under $Q$ this follows from the first subspace moment; under $P$, the
likelihood identity bounds it by a constant times the first two subspace
moments under $Q$, using $m_K^2=m_K+q(q+1)\qbinom K2$.
The nullspace computation costs $O(n^6)$ field
operations and each determinant test costs $O(n^3)$, so the complete
optimal procedure has expected $O_{q,a}(n^6)$ field-operation cost.
The limiting success probability of elimination alone is $g_\infty/g_a$
for $a\ge0$ and zero for $a<0$, since its success event under the planted
law is $\dim\mathcal N=1$. This expectation statement concerns the
specified random experiments; the worst-case nullspace can be large.

\paragraph{Numerical example.}
For $q=2$, at the parameter-count threshold $t=D-1$ ($a=0$),
\[
 \begin{array}{ll}
 T_2(0)=0.411496025790588369\ldots,&
 T_2^{\rm dim}(0)=0.288788095086602421\ldots,\\
 S_2(0)=0.701564718009178670\ldots,&
 U_2(0)=0.488389858128608926\ldots.
 \end{array}
\]
Elimination alone recovers a unique linear form with limiting probability
$\prod_{j\ge1}(1-2^{-j})=0.288788095086602421\ldots$. Checking nondegeneracy resolves
additional cases, and even when several compatible nondegenerate forms
remain, choosing one uniformly achieves the optimal success $S_2(0)$.
